\section{Battery Witness Results}
\label{sec:prof-battery-results}

The battery domain remains the strongest empirical witness in ORIUS, but it is
best read as a witness rather than the conceptual center of the framework.  The
locked battery artifacts show a clear separation between \emph{controllers that
can detect and repair under degraded observation} and those that cannot.  That
separation matters more than any single cost number: the decisive result is
whether true-state violations disappear under telemetry faults.

Table~\ref{tab:prof-battery-summary} condenses the strongest locked battery
evidence.  The repair-aware ORIUS variants are the only methods in this slice
that simultaneously eliminate true-state violations and retain meaningful
coverage.  The deterministic LP and static interval baselines fail on that
first criterion, while the wrapped and FTIT variants preserve safety at the
price of modest intervention.

\begin{table}[t]
\centering
\small
\caption{Compact battery witness summary from locked publication artifacts.
TSVR is true-state violation rate; coverage is the 90\% predictive interval
coverage reported by the baseline audit.}
\label{tab:prof-battery-summary}
\begin{tabular}{lrrrr}
\toprule
\textbf{Controller} & \textbf{TSVR (\%)} & \textbf{Repair / IR (\%)} &
\textbf{Coverage} & \textbf{Mean width} \\
\midrule
ORIUS FTIT & 0.000 & 2.798 & 0.971627 & 3906.548 \\
ORIUS Wrapped & 0.000 & 0.000 & 0.884127 & 559.921 \\
Adaptive Conformal + Repair & 0.000 & 1.389 & 0.962500 & 2824.252 \\
Deterministic LP & 3.909 & 0.000 & 0.883333 & 503.131 \\
Robust Fixed Interval & 25.575 & 0.000 & 0.883333 & 503.131 \\
\bottomrule
\end{tabular}
\end{table}

Two points follow directly from the evidence.  First, the safety gain is
structural: once reliability-aware repair is enabled, the hidden violation rate
drops to zero across the battery witness package.  Second, the deep quality
model is not being used as a claim of forecasting superiority.  The deep OQE
summary is only mildly predictive of fault family (\texttt{fault\_auc} $=0.534$),
which is precisely why the framework treats quality scoring as a guardrail for
repair and certificate logic rather than as a standalone predictor.

Figure~\ref{fig:prof-battery-baselines} provides the visual counterpart to the
table: the battery witness is strongest where the repair-and-certify layer keeps
the true-state violation surface at zero while leaving the non-repair baselines
exposed.

\begin{figure}[t]
\centering
\includegraphics[width=\linewidth]{fig_battery_reliability_baselines}
\caption{Battery reliability baselines from the locked publication package.
The ORIUS repair-aware variants preserve zero true-state violations, while the
deterministic and static-interval baselines do not.}
\label{fig:prof-battery-baselines}
\end{figure}

The chapter-level battery results remain consistent with the publication
artifacts used elsewhere in the monograph: the locked controller comparison,
the reliability baseline audit, the 48-hour trace, and the ablation suite all
point to the same conclusion.  ORIUS is not merely reducing cost or widening
intervals; it is preventing hidden violations under observation degradation.

\paragraph{Bottom line.}
For a professor-facing ORIUS manuscript, this is the battery story worth
retaining in the main text: the battery plant is the cleanest empirical witness
that the observation-quality signal, repair step, and certificate discipline are
doing real work.  The detailed traces, fault-family responses, and ablations are
valuable, but they belong behind the compact result statement rather than
crowding it.
